\section{Analytic/period column}
\label{sec:analytic-period}

The analytic/period column keeps the analytic leading coefficient and the
real period apart from the height, finite and local data.  Its native data
form a pair of numbers, and its readout is their quotient.

\begin{definition}[Real analytic/period map]\label{def:apparatus:analytic-period-map}
By modularity~\citep{Wiles1995,BreuilConradDiamondTaylor2001},
$L(E,s)=L(f_E,s)$ extends to an entire function.  Let
$r_{\mathrm{an}}=\operatorname{ord}_{s=1}L(E,s)$ and
\[
  A_E=\frac{L^{(r_{\mathrm{an}})}(E,1)}{r_{\mathrm{an}}!}\neq0
\]
be the analytic leading coefficient, that is, the leading Taylor
coefficient of $L(E,s)$ at $s=1$.  Let $\Omega_E^+>0$ be the real period
in the BSD normalization: the integral of $|\omega_E|$ over $E(\R)$ for a
N\'eron differential $\omega_E$.  The native analytic/period family is the
pair
\[
  L_{\mathrm{an/per}}(E)=(A_E,\Omega_E^+),
\]
and the real analytic/period map is
\[
  \rho_{\mathrm{an/period}}^{\R}:(A_E,\Omega_E^+)\longmapsto
  (A_E/\Omega_E^+)\,e_{\mathrm{an/per}},
\]
where $e_{\mathrm{an/per}}$ is a basis vector of a real line.  The column
contains no height, torsion, Tamagawa or $\Sha$ data.
\end{definition}

The adequacy test asks whether the pair is enough to recover the first-order
variation of the quotient.  In logarithmic coordinates the absolute value
of the quotient becomes a difference; its sign is not a tangent direction
and is recorded separately, as part of the orientation data of
\Cref{sec:det-assembly}.

\begin{theorem}[Analytic/period tangent Schur collapse]\label{thm:apparatus:analytic-period-schur-collapse}
In the chart
\[
  x=(x_A,x_\Omega)=(\log|A_E|,\log\Omega_E^+)\in\R^2,
\]
take $C=I_2$, the full native source $L_{\mathrm{full}}(x)=x$, and the
readout $D(x)=x_A-x_\Omega=\log|A_E/\Omega_E^+|=dx$ with $d=(1,-1)$.  Then $\KLL=I_2$,
$\KDD=dd^{\mathsf T}=2$ and
\[
  \XiC=2-d\,I_2\,d^{\mathsf T}=0.
\]
With no source admitted the residual is $2$; with only $x_A$ admitted it
is $1$; with both coordinates admitted it is $0$.%
\footnote{Lean proves these three integer residual computations for the
row $(1,-1)$.}
\end{theorem}

\begin{proof}
The readout energy is $\KDD=1^2+(-1)^2=2$.  With the full source,
$\KLL=\Kdagger=I_2$, $\KDL=d$ and $K_{LD}=d^{\mathsf T}$, so the projected
energy is $d\,I_2\,d^{\mathsf T}=2$ and $\XiC=0$.  With no source nothing is
projected away and the residual is $2$.  If only $x_A$ is admitted, $d$
splits orthogonally as the seen part $(1,0)$ plus the unseen part
$(0,-1)$, and the residual is the squared length $1$ of the unseen part.
\end{proof}

The theorem is elementary, and that is the point: a single observed
coordinate cannot recover a quotient of two independent quantities,
while the pair can.  The arithmetic content of the column lies entirely in
the comparison below.

\begin{obligation}[Analytic/period transport]\label{obl:apparatus:rho-an-period-transport}
The Bloch--Kato analytic component requires a comparison
\[
  \beta_{\mathrm{BK,an}}:
  L_{\mathrm{an/per}}(E)_{\R}
  \longrightarrow
  \Delta_f(M_E)_{\mathrm{an/per}}
\]
identifying $A_E/\Omega_E^+$ with the zeta-element factor of the
fundamental line.  It must fix the rational and integral normalizations,
the Betti--de Rham comparison and the real orientation, the leading-term
convention in positive rank, and the separation from the other four
columns.  We treat this comparison as external input in every rank.  In
rank $0$ the analytic quotient $L(E,1)/\Omega_E^+$ is a classical object,
but matching it to the zeta-element factor in the conventions above is
part of the comparison, not a consequence of the Schur calculation.
\end{obligation}
